\paragraph{\texttt{np.arange(start, stop)}} Integers from start up to stop-1. We use it for the grid index i = 1,...,7.

\begin{lstlisting}[style=py]
import numpy as np
i = np.arange(1, 8)
print(i)
print(i/8)
\end{lstlisting}

\begin{lstlisting}[style=out]
[1 2 3 4 5 6 7]
[0.125 0.25  0.375 0.5   0.625 0.75  0.875]
\end{lstlisting}

\paragraph{\texttt{np.ones(n) and np.zeros((m, n))}} Arrays filled with ones or zeros. zeros((m, n)) reserves space for the snapshot matrix.

\begin{lstlisting}[style=py]
import numpy as np
print(np.ones(5))
print(np.zeros((2, 5)))
\end{lstlisting}

\begin{lstlisting}[style=out]
[1. 1. 1. 1. 1.]
[[0. 0. 0. 0. 0.]
 [0. 0. 0. 0. 0.]]
\end{lstlisting}

\paragraph{\texttt{np.diag(v, k)}} Puts the vector v on the k-th diagonal of a square matrix (k = 0 main, k = 1 above, k = -1 below).

\begin{lstlisting}[style=py]
import numpy as np
print(np.diag([-2, -2, -2, -2, -2]))
print(np.diag([1, 1, 1, 1], 1))
\end{lstlisting}

\begin{lstlisting}[style=out]
[[-2  0  0  0  0]
 [ 0 -2  0  0  0]
 [ 0  0 -2  0  0]
 [ 0  0  0 -2  0]
 [ 0  0  0  0 -2]]
[[0 1 0 0 0]
 [0 0 1 0 0]
 [0 0 0 1 0]
 [0 0 0 0 1]
 [0 0 0 0 0]]
\end{lstlisting}

\paragraph{\texttt{Building the 1D Laplacian}} The sum of the three diagonals gives the tridiagonal matrix of Section 6 (here N = 5, without the 1/dx**2 factor).

\begin{lstlisting}[style=py]
import numpy as np
N = 5
T = np.diag(-2*np.ones(N)) + np.diag(np.ones(N-1), 1) + np.diag(np.ones(N-1), -1)
print(T)
\end{lstlisting}

\begin{lstlisting}[style=out]
[[-2.  1.  0.  0.  0.]
 [ 1. -2.  1.  0.  0.]
 [ 0.  1. -2.  1.  0.]
 [ 0.  0.  1. -2.  1.]
 [ 0.  0.  0.  1. -2.]]
\end{lstlisting}

\paragraph{\texttt{np.eye(n)}} The identity matrix. A = I + dt L is one explicit Euler step.

\begin{lstlisting}[style=py]
import numpy as np
print(np.eye(5))
\end{lstlisting}

\begin{lstlisting}[style=out]
[[1. 0. 0. 0. 0.]
 [0. 1. 0. 0. 0.]
 [0. 0. 1. 0. 0.]
 [0. 0. 0. 1. 0.]
 [0. 0. 0. 0. 1.]]
\end{lstlisting}

\paragraph{\texttt{The @ operator}} Matrix times vector (or matrix times matrix). One FOM step is A @ x.

\begin{lstlisting}[style=py]
import numpy as np
T = np.diag(-2*np.ones(5)) + np.diag(np.ones(4), 1) + np.diag(np.ones(4), -1)
x = np.array([0., 1., 2., 1., 0.])
print(T @ x)
\end{lstlisting}

\begin{lstlisting}[style=out]
[ 1.  0. -2.  0.  1.]
\end{lstlisting}

\paragraph{\texttt{Filling columns: X[:, k]}} X[:, k] is column k of X. The time loop writes the state x(k) into column k.

\begin{lstlisting}[style=py]
import numpy as np
X = np.zeros((3, 4))
X[:, 0] = [1., 2., 3.]
for k in range(3):
    X[:, k+1] = 0.5 * X[:, k]
print(X)
\end{lstlisting}

\begin{lstlisting}[style=out]
[[1.    0.5   0.25  0.125]
 [2.    1.    0.5   0.25 ]
 [3.    1.5   0.75  0.375]]
\end{lstlisting}

\paragraph{\texttt{Slicing: M[:, :-1] and M[:, 1:]}} All columns except the last, and all columns except the first. Side by side they pair x(k) with x(k+1).

\begin{lstlisting}[style=py]
import numpy as np
M = np.arange(10).reshape(2, 5)
print(M)
print(M[:, :-1])
print(M[:, 1:])
\end{lstlisting}

\begin{lstlisting}[style=out]
[[0 1 2 3 4]
 [5 6 7 8 9]]
[[0 1 2 3]
 [5 6 7 8]]
[[1 2 3 4]
 [6 7 8 9]]
\end{lstlisting}

\paragraph{\texttt{Transpose .T}} Swaps rows and columns. OpInf stacks the reduced states as ROWS of D, so we transpose.

\begin{lstlisting}[style=py]
import numpy as np
M = np.arange(6).reshape(2, 3)
print(M)
print(M.T)
\end{lstlisting}

\begin{lstlisting}[style=out]
[[0 1 2]
 [3 4 5]]
[[0 3]
 [1 4]
 [2 5]]
\end{lstlisting}

\paragraph{\texttt{np.linalg.svd(X, full\_matrices=False)}} X = U diag(s) Wt. Columns of U are the POD modes, s are the singular values (largest first). A rank-1 example: only one singular value is nonzero.

\begin{lstlisting}[style=py]
import numpy as np
np.set_printoptions(precision=4, suppress=True)
X = np.outer([1., 2., 2.], [1., 0.5, 0.25, 0.125, 0.0625])
U, s, Wt = np.linalg.svd(X, full_matrices=False)
print(s)
print(U[:, 0])
\end{lstlisting}

\begin{lstlisting}[style=out]
[3.4624 0.     0.    ]
[-0.3333 -0.6667 -0.6667]
\end{lstlisting}

\paragraph{\texttt{U[:, :r]}} The first r columns: the POD basis V.

\begin{lstlisting}[style=py]
import numpy as np
U = np.arange(15).reshape(5, 3)
print(U[:, :2])
\end{lstlisting}

\begin{lstlisting}[style=out]
[[ 0  1]
 [ 3  4]
 [ 6  7]
 [ 9 10]
 [12 13]]
\end{lstlisting}

\paragraph{\texttt{np.linalg.lstsq(D, R, rcond=None)}} Finds the M that makes the residual D M - R as small as possible (least squares). Here the data come from x(k+1) = 0.8 x(k) exactly, so the fitted slope is 0.8.

\begin{lstlisting}[style=py]
import numpy as np
x = 0.8**np.arange(6)            # 1, 0.8, 0.64, ...
D = x[:-1].reshape(-1, 1)        # column of x_k
R = x[1:].reshape(-1, 1)         # column of x_{k+1}
M, res, rank, sv = np.linalg.lstsq(D, R, rcond=None)
print(M)
\end{lstlisting}

\begin{lstlisting}[style=out]
[[0.8]]
\end{lstlisting}

\paragraph{\texttt{np.linalg.norm and np.linalg.eigvals}} The (Frobenius) norm measures the size of a difference; eigvals gives the eigenvalues of a matrix.

\begin{lstlisting}[style=py]
import numpy as np
np.set_printoptions(precision=4, suppress=True)
B = np.array([[0.9, 0.1], [0.1, 0.7]])
print(np.linalg.norm(B - np.diag([0.9, 0.7])))
print(np.linalg.eigvals(B))
\end{lstlisting}

\begin{lstlisting}[style=out]
0.14142135623730953
[0.9414+0.j 0.6586+0.j]
\end{lstlisting}
